\documentclass[11pt]{article}
\usepackage[a4paper, top=18mm, bottom=18mm, left=20mm, right=20mm]{geometry}
\usepackage[T2A]{fontenc}
\usepackage[utf8]{inputenc}
\usepackage[ukrainian]{babel}
\usepackage{graphicx}
\usepackage[export]{adjustbox}
\usepackage{amsmath}
\usepackage{amssymb}
\graphicspath{{/app/Quiz/Scripts/2023/ProgAll/15BinTree/}{/app/Quiz/Scripts/2021/Mod2021_3/BinTree/}{/app/Quiz/Scripts/2020/Mod2020_4/BinTree/}}
\setlength{\parindent}{0pt}
\setlength{\parskip}{6pt}
\pagestyle{empty}
\begin{document}
%begin
{\large\textbf{Контрольна робота}}

\textbf{Варіант \No\,1}\hfill \textit{ПІБ:}\,\underline{\hspace{60mm}}
\vspace{4mm}

%beginitem
1. %goodpath:Scripts/2018/Mod2018_1/03-good.txt
Відмітити все, що є коректним оператором С++:
( 1 ) cout<<(3=!5);

( 2 ) bool f(int a);

( 3 ) cin<<x;

( 4 ) j=i;
\vspace{5mm}
%enditem
%beginitem
2. Що буде виведено за виконання фрагменту коду?

%code
4**2.0+2
%code
\vspace{5mm}
%enditem
%beginitem
3. Що буде виведено за виконання фрагменту коду?

%code
print(5.0 > 3.0 <= True != 5)
%code
\vspace{5mm}
%enditem
%beginitem
4. %goodpath:Scripts/2019/Mod2019_1/Lex/01-good.txt
Відмітити все, що є однією правильною лексемою:
( 1 ) (1)

( 2 ) 0.2E3

( 3 ) "abc'

( 4 ) x23
\vspace{5mm}
%enditem
%beginitem
5. Що буде виведено за виконання фрагменту коду?

a,b,c=3,1,5
def h(a):
global c    a=3
    b*=2
    c=4
    return a+b+c

a,b,c=0,9,6
print(h(b),a,b,c)
\vspace{5mm}
%enditem
%beginitem
6. Що буде виведено за виконання фрагменту коду?

%code
cout << (5 % 5 % 5 % 5);
%code
\vspace{5mm}
%enditem
%beginitem
7. 

\begin{center}
	\adjustimage{max size={0.8\linewidth}{0.8\paperheight}}
	{../15BinTree/trees_exp/59.png}
\end{center}
Указати послідовність міток вершин дерева, зображеного на рис., що утвориться в результаті обходу дерева в прямому порядку. Мітки записувати через пробіл.\par Алгоритм починає роботу з кореня (на рис. це верхня вершина).
%answer:59 прямому
\vspace{5mm}
%enditem
%beginitem
8. %goodpath:Scripts/2019/Mod2019_1/Lex/03-good.txt
Відмітити все, що є коректним оператором С++:
( 1 ) if a<b: a=0 else: b=0

( 2 ) x=3**1

( 3 ) if ('x'><'y'): a=6

( 4 ) --n
\vspace{5mm}
%enditem
%beginitem
9. Що буде виведено за виконання фрагменту коду?
%code
a, b, c = 9, 7, 8
def h(a, b=6, c=7):
    print(a, b, c, end=" ")

h(3, 1, c=5)
print(a, b, c)
%code
\vspace{5mm}
%enditem
%beginitem
10. Що буде виведено за виконання фрагменту коду?

print('R', end=' ')
d=[10,11,12,13,14]
print(d.pop(-2), end=' ')
print(d)

\vspace{5mm}
%enditem










%end
\newpage

\newpage

%begin
{\large\textbf{Контрольна робота}}

\textbf{Варіант \No\,2}\hfill \textit{ПІБ:}\,\underline{\hspace{60mm}}
\vspace{4mm}

%beginitem
1. Що буде виведено за виконання фрагменту коду?

a,b,c=3,1,5
def h(a):
global c    a=3
    b*=2
    c=4
    return a+b+c

a,b,c=0,9,6
print(h(b),a,b,c)
\vspace{5mm}
%enditem
%beginitem
2. Що буде виведено за виконання фрагменту коду?

%code
print('R', end=' ')
d = ('a','b','c','d','e',0,1,2,3)
print(d[-11])
%code

\vspace{5mm}
%enditem
%beginitem
3. %goodpath:Scripts/2018/Mod2018_1/03-good.txt
Відмітити все, що є коректним оператором С++:
( 1 ) cout<<(3=!5);

( 2 ) bool f(int a);

( 3 ) cin<<x;

( 4 ) j=i;
\vspace{5mm}
%enditem
%beginitem
4. Що буде виведено за виконання фрагменту коду?

%code
cout << (!5.0>=3.0 || !true<= 5 || 3!=7);
%code
\vspace{5mm}
%enditem
%beginitem
5. Що буде виведено за виконання фрагменту коду?
%code

a,b,c=9,7,8
def h(a,b=6,c=7):
    print(a,b,c,end="")

h(3,1,c=5)
print(a,b,c)

%code
\vspace{5mm}
%enditem
%beginitem
6. Що буде виведено за виконання фрагменту коду?

x,y,z=3,1,5
x,y,z,x=9,z,y,7
print(x,y,z)
\vspace{5mm}
%enditem
%beginitem
7. %goodpath:Scripts/2019/Mod2019_1/Lex/03-good.txt
Відмітити все, що є коректним оператором С++:
( 1 ) if a<b: a=0 else: b=0

( 2 ) x=3**1

( 3 ) if ('x'><'y'): a=6

( 4 ) --n
\vspace{5mm}
%enditem
%beginitem
8. %goodpath:Scripts/2018/Mod2018_1/01-good.txt
Відмітити все, що є однією правильною лексемою:
( 1 ) 0,5

( 2 ) LIFO

( 3 ) x+23

( 4 ) x23
\vspace{5mm}
%enditem
%beginitem
9. Що буде виведено за виконання фрагменту коду?

%code
print(4**2.0+2)
%code
\vspace{5mm}
%enditem
%beginitem
10. Що буде виведено за виконання фрагменту коду?

%code
#include <iostream>
using namespace std;

int h(int d){
    int x = 69;
    if (d >= -3) 
        return 3;
    if (d <= -4)
         x = 1;
    else
         return 5;
    return x;
}

int main(){
    cout << h(9);
    return 0;
}
%code

\vspace{5mm}
%enditem










%end
\newpage

\newpage

%begin
{\large\textbf{Контрольна робота}}

\textbf{Варіант \No\,3}\hfill \textit{ПІБ:}\,\underline{\hspace{60mm}}
\vspace{4mm}

%beginitem
1. Що буде виведено за виконання фрагменту коду?

%code
cout << 5 % 5 % 5 % 5;
%code
\vspace{5mm}
%enditem
%beginitem
2. %goodpath:Scripts/2019/Mod2019_1/Lex/01-good.txt
Відмітити все, що є однією правильною лексемою:
( 1 ) (1)

( 2 ) 0.2E3

( 3 ) "abc'

( 4 ) x23
\vspace{5mm}
%enditem
%beginitem
3. Що буде виведено за виконання фрагменту коду?
try:
    res = 4**2.0+2
    if isinstance(res, bool):
        print(f'bool;{res}')
    elif isinstance(res, int):
        print(f'int;{res}')
    elif isinstance(res, float):
        print(f'float;{res:.2f}')
    else:
        print('???')
except:
    print('Z')
\vspace{5mm}
%enditem
%beginitem
4. Що буде виведено за виконання фрагменту коду?
try:
    res = 5/5%5*5
    if isinstance(res, bool):
        print(f'bool;{res}')
    elif isinstance(res, int):
        print(f'int;{res}')
    elif isinstance(res, float):
        print(f'float;{res:.2f}')
    else:
        print('???')
except:
    print('Z')
\vspace{5mm}
%enditem
%beginitem
5. 
Вкажіть значення та тип виразу:
"3">=3.0
\vspace{5mm}
%enditem
%beginitem
6. %goodpath:Scripts/2019/Mod2019_1/Lex/03-good.txt
Відмітити все, що є коректним оператором С++:
( 1 ) if a<b: a=0 else: b=0

( 2 ) x=3**1

( 3 ) if ('x'><'y'): a=6

( 4 ) --n
\vspace{5mm}
%enditem
%beginitem
7. Що буде виведено за виконання фрагменту коду?
try:
    for d in range(3, 3, -3):
        if d>=6:
            break
        print(d, end=';')
        if d>=6:
            break
    else:
        print(d,end=';')
    print(d)
except:
    print('Z')
\vspace{5mm}
%enditem
%beginitem
8. Що буде виведено за виконання фрагменту коду?
%code
#include <iostream>
using namespace std;

int f(int &x, int &y){
    x = 3;
    y+= 2;
    return x;
}

int main(){
    int a = 6, b = 5;
    a = f(a, a);
    cout << a << ":" << b <<':';
    {
        int a = 1, b = 4;
        cout << ((a>=5) && ((b+=1) >= 5))
             << a << ":" << b << ":";
    }
    {
        inta = 8;
        cout << a << ':';
    }
    cout << a << ":" << b << endl;
    return 0;
}
%code

\vspace{5mm}
%enditem
%beginitem
9. %goodpath:Scripts/2018/Mod2018_1/01-good.txt
Відмітити все, що є однією правильною лексемою:
( 1 ) 0,5

( 2 ) LIFO

( 3 ) x+23

( 4 ) x23
\vspace{5mm}
%enditem
%beginitem
10. Що буде виведено за виконання фрагменту коду?

%code
print('R', end=' ')
n=3
while n<7:
    n+=1
print(n)
%code

\vspace{5mm}
%enditem










%end
\newpage

\newpage

%begin
{\large\textbf{Контрольна робота}}

\textbf{Варіант \No\,4}\hfill \textit{ПІБ:}\,\underline{\hspace{60mm}}
\vspace{4mm}

%beginitem
1. Що буде виведено за виконання фрагменту коду?
a = 3
b = 1
c = 5

def h(a):
global c    a = 3
    b *= 2
    c = 4
    return a + b + c

a = 0
b = 9
c = 6

try:
    print(h(b), a, b, c, sep=';')
except:
    print('Z', a, b, c, sep=';')
\vspace{5mm}
%enditem
%beginitem
2. 
Записати ВИРАЗ, значенням якого є множина, що складається з цілих чисел від 12 до 2018 (включно), що не діляться на 7.\vspace{5mm}
%enditem
%beginitem
3. %goodpath:Scripts/2019/Mod2019_1/Lex/03-good.txt
Відмітити все, що є коректним оператором С++:
( 1 ) if a<b: a=0 else: b=0

( 2 ) x=3**1

( 3 ) if ('x'><'y'): a=6

( 4 ) --n
\vspace{5mm}
%enditem
%beginitem
4. 
Записати ВИРАЗ, значенням якого є список, що складається з кубів цілих чисел від 12 до 2018 (включно), що не діляться на жодне з чисел 2, 3, записаних в оберненому порядку.\vspace{5mm}
%enditem
%beginitem
5. %goodpath:Scripts/2018/Mod2018_1/01-good.txt
Відмітити все, що є однією правильною лексемою:
( 1 ) 0,5

( 2 ) LIFO

( 3 ) x+23

( 4 ) x23
\vspace{5mm}
%enditem
%beginitem
6. 
Написати програму, що запитує у користувача значення дійсних параметрів $x$ та $y$, обчислює для них значення виразу 
$$\frac{\log_{2} x + 12}{(\tg x + 0.3) (y + 68)}$$ 
та повiдомляє введенi данi та результат обчислення.
 Оцінюється правильність та якість коду.

\vspace{5mm}
%enditem
%beginitem
7. %goodpath:Scripts/2019/Mod2019_1/Lex/01-good.txt
Відмітити все, що є однією правильною лексемою:
( 1 ) (1)

( 2 ) 0.2E3

( 3 ) "abc'

( 4 ) x23
\vspace{5mm}
%enditem
%beginitem
8. Що буде виведено за виконання фрагменту коду?

print('R', end=' ')
for d in range(3,3,-3):
    if d>=6:
        break
        print(d, end=' ')
    if d>=6:
        break
    else:
        print(d, end=' ')
    print(d)
\vspace{5mm}
%enditem
%beginitem
9. Що буде виведено за виконання фрагменту коду?
d=[10,11,12,13,14]; print(d.pop(-2),end=' '); print(d)\vspace{5mm}
%enditem
%beginitem
10. 

\begin{center}
	\adjustimage{max size={0.8\linewidth}{0.8\paperheight}}
	{../15BinTree/trees_exp/59.png}
\end{center}
Указати послідовність міток вершин дерева, зображеного на рис., що утвориться в результаті обходу дерева в прямому порядку. Мітки записувати через пробіл.\par Алгоритм починає роботу з кореня (на рис. це верхня вершина).
%answer:59 прямому
\vspace{5mm}
%enditem










%end
\newpage

\newpage

%begin
{\large\textbf{Контрольна робота}}

\textbf{Варіант \No\,5}\hfill \textit{ПІБ:}\,\underline{\hspace{60mm}}
\vspace{4mm}

%beginitem
1. Що буде виведено за виконання фрагменту коду?

%code
print(4**2.0+2)
%code
\vspace{5mm}
%enditem
%beginitem
2. %goodpath:Scripts/2019/Mod2019_1/Lex/01-good.txt
Відмітити все, що є однією правильною лексемою:
( 1 ) (1)

( 2 ) 0.2E3

( 3 ) "abc'

( 4 ) x23
\vspace{5mm}
%enditem
%beginitem
3. 

\begin{center}
	\adjustimage{max size={0.8\linewidth}{0.8\paperheight}}
	{../BinTree/trees_exp/59.png}
\end{center}
Указати послідовність міток вершин дерева, зображеного на рис., що утвориться в результаті обходу дерева в прямому порядку. Мітки записувати через пробіл.\par Алгоритм починає роботу з кореня (на рис. це верхня вершина).
%answer:59 прямому
\vspace{5mm}
%enditem
%beginitem
4. Що буде виведено за виконання фрагменту коду?

print('R', end=' ')
n=3
while n<7:
    n+=1
print(n)

\vspace{5mm}
%enditem
%beginitem
5. 
Записати ВИРАЗ, значенням якого є множина, що складається з кубів цілих чисел від 12 до 2018 (включно), що не діляться на жодне з чисел 2, 3.\vspace{5mm}
%enditem
%beginitem
6. %goodpath:Scripts/2019/Mod2019_1/Lex/03-good.txt
Відмітити все, що є коректним оператором С++:
( 1 ) if a<b: a=0 else: b=0

( 2 ) x=3**1

( 3 ) if ('x'><'y'): a=6

( 4 ) --n
\vspace{5mm}
%enditem
%beginitem
7. Що буде виведено за виконання фрагменту коду?
try:
    n = 3
    while n<7:
        n += 1
    print(n, end=';')
except:
    print('ERR')
\vspace{5mm}
%enditem
%beginitem
8. %goodpath:Scripts/2018/Mod2018_1/01-good.txt
Відмітити все, що є однією правильною лексемою:
( 1 ) 0,5

( 2 ) LIFO

( 3 ) x+23

( 4 ) x23
\vspace{5mm}
%enditem
%beginitem
9. Що буде виведено за виконання фрагменту коду?

%code
#include <iostream>
using namespace std;

int a = 3, b = 1, c = 5;

int h(){
    inta = 3;
    intb = 0;
    intc = 9;
    return a + b + c;
}

int main(){
    inta = 6;
    intb = 4;
    intc = 8;
    cout << h() << ':';
    cout << a << ':' << b << ':' << c;
    return 0;
}
%code
\vspace{5mm}
%enditem
%beginitem
10. Що буде виведено за виконання фрагменту коду?

%code
a,b,c=3,1,5
def h(a):
global c    a=3
    b*=2
    c=4
    return a+b+c

a,b,c=0,9,6
print(h(b),a,b,c)
%code
\vspace{5mm}
%enditem










%end
\newpage

\newpage

%begin
{\large\textbf{Контрольна робота}}

\textbf{Варіант \No\,6}\hfill \textit{ПІБ:}\,\underline{\hspace{60mm}}
\vspace{4mm}

%beginitem
1. Що буде виведено за виконання фрагменту коду?

%code
a,b,c=3,1,5
def h(a):
global c    a=3
    b*=2
    c=4
    return a+b+c

a,b,c=0,9,6
print(h(b),a,b,c)
%code
\vspace{5mm}
%enditem
%beginitem
2. %goodpath:Scripts/2019/Mod2019_1/Lex/03-good.txt
Відмітити все, що є коректним оператором С++:
( 1 ) if a<b: a=0 else: b=0

( 2 ) x=3**1

( 3 ) if ('x'><'y'): a=6

( 4 ) --n
\vspace{5mm}
%enditem
%beginitem
3. Що буде виведено за виконання фрагменту коду?

%code
for d in range(8,8+4,3):
    if d>=6:
        pass
    print(d, end=' ')
    d=8
    if d>=6:
        break
else:
    print(d, end=' ')
print(d, end=' ')
%code
\vspace{5mm}
%enditem
%beginitem
4. Що буде виведено за виконання фрагменту коду?

%code
print(5.0>3.0>True>5)
%code
\vspace{5mm}
%enditem
%beginitem
5. Що буде виведено за виконання фрагменту коду?

%code
a = 3
b = 1
c = 5

def h(a):
global c    a=3
    b=2
    c=4
    return a+b+c

a, b, c = 0,9,6
try:
    print(h(b), a, b, c, sep=';')
except:
    print('Z', a, b, c, sep=';')
%code
\vspace{5mm}
%enditem
%beginitem
6. Що буде виведено за виконання фрагменту коду?

%code
print('R', end=' ')
n=7
while n>3:
    n+=-3
print(n)
%code

\vspace{5mm}
%enditem
%beginitem
7. Що буде виведено за виконання фрагменту коду?

%code
def h(d):
    x=69
    if d>=-3: 
        return 3
    if d<=-4:
         x=1
    else:
         return 5
    return x

print(h(9))
%code
\vspace{5mm}
%enditem
%beginitem
8. Що буде виведено за виконання фрагменту коду?

%code
not5.0>=3.0 and notTrue<=5 and 3!=7
%code
\vspace{5mm}
%enditem
%beginitem
9. %goodpath:Scripts/2018/Mod2018_1/03-good.txt
Відмітити все, що є коректним оператором С++:
( 1 ) cout<<(3=!5);

( 2 ) bool f(int a);

( 3 ) cin<<x;

( 4 ) j=i;
\vspace{5mm}
%enditem
%beginitem
10. %goodpath:Scripts/2019/Mod2019_1/Lex/01-good.txt
Відмітити все, що є однією правильною лексемою:
( 1 ) (1)

( 2 ) 0.2E3

( 3 ) "abc'

( 4 ) x23
\vspace{5mm}
%enditem










%end
\newpage

\newpage

%begin
{\large\textbf{Контрольна робота}}

\textbf{Варіант \No\,7}\hfill \textit{ПІБ:}\,\underline{\hspace{60mm}}
\vspace{4mm}

%beginitem
1. Що буде виведено за виконання фрагменту коду?

def h(d):
    x=69
    if d>=-3: 
        return 3
    if d<=-4:
         x=1
    else:
         return 5
    return x

print(h(9))
\vspace{5mm}
%enditem
%beginitem
2. %goodpath:Scripts/2019/Mod2019_1/Lex/03-good.txt
Відмітити все, що є коректним оператором С++:
( 1 ) if a<b: a=0 else: b=0

( 2 ) x=3**1

( 3 ) if ('x'><'y'): a=6

( 4 ) --n
\vspace{5mm}
%enditem
%beginitem
3. Що буде виведено за виконання фрагменту коду?

%code
a,b,c=3,1,5
def h(a):
    a=3
    b=2
    c=4
    return a+b+c

a,b,c=0,9,6
print(h(b),a,b,c)
%code
\vspace{5mm}
%enditem
%beginitem
4. 
Написати програму, що запитує у користувача значення дійсних параметрів $x$ та $y$, обчислює для них значення виразу та повiдомляє введенi данi та результат обчислення.
$$\frac{\log_{2} x + 12}{(\tg x + 0.3) (y + 68)}$$ 

\vspace{5mm}
%enditem
%beginitem
5. %goodpath:Scripts/2018/Mod2018_1/01-good.txt
Відмітити все, що є однією правильною лексемою:
( 1 ) 0,5

( 2 ) LIFO

( 3 ) x+23

( 4 ) x23
\vspace{5mm}
%enditem
%beginitem
6. Що буде виведено за виконання фрагменту коду?

%code
n=3
while n<7:
    n+=1
print(n)
%code

\vspace{5mm}
%enditem
%beginitem
7. Що буде виведено за виконання фрагменту коду?

%code
def h(a,b):
    c=69
    if b>=-3:
        return 3
    if a<=-4:
         c=1
    else: 
        return 5
    return c

print(h(9,9))
%code
\vspace{5mm}
%enditem
%beginitem
8. %goodpath:Scripts/2018/Mod2018_1/03-good.txt
Відмітити все, що є коректним оператором С++:
( 1 ) cout<<(3=!5);

( 2 ) bool f(int a);

( 3 ) cin<<x;

( 4 ) j=i;
\vspace{5mm}
%enditem
%beginitem
9. Що буде виведено за виконання фрагменту коду?

def f(n):
    print("f", end="");
    return n

print(f(9) and f(-8))
\vspace{5mm}
%enditem
%beginitem
10. Що буде виведено за виконання фрагменту коду?

%code
cout << (5 % 5 % 5 % 5);
%code
\vspace{5mm}
%enditem










%end
\newpage

\newpage

%begin
{\large\textbf{Контрольна робота}}

\textbf{Варіант \No\,8}\hfill \textit{ПІБ:}\,\underline{\hspace{60mm}}
\vspace{4mm}

%beginitem
1. Що буде виведено за виконання фрагменту коду?
%code
a, b, c = 9, 7, 8
def h(a, b=6, c=7):
    print(a, b, c, end=" ")

h(3, 1, c=5)
print(a, b, c)
%code
\vspace{5mm}
%enditem
%beginitem
2. Що буде виведено за виконання фрагменту коду?

%code
n=3
while n<7:
    n+=1
print(n)
%code

\vspace{5mm}
%enditem
%beginitem
3. Що буде виведено за виконання фрагменту коду?

def f(n):
    print("f", end="");
    return n

print(f(9) and f(-8))
\vspace{5mm}
%enditem
%beginitem
4. %goodpath:Scripts/2019/Mod2019_1/Lex/01-good.txt
Відмітити все, що є однією правильною лексемою:
( 1 ) (1)

( 2 ) 0.2E3

( 3 ) "abc'

( 4 ) x23
\vspace{5mm}
%enditem
%beginitem
5. Що буде виведено за виконання фрагменту коду?
%code
#include <iostream>
using namespace std;

int f(int &x, int &y){
    x = 3;
    y+= 2;
    return x;
}

int main(){
    int a = 6, b = 5;
    a = f(a, a);
    cout << a << ":" << b <<':';
    {
        int a = 1, b = 4;
        cout << ((a>=5) && ((b+=1) >= 5)) << ':';
        cout << a << ":" << b << ':';
    }
    {
        inta = 8;
        cout << a << ':';
    }
    cout << a << ":" << b << endl;
    return 0;
}
%code

\vspace{5mm}
%enditem
%beginitem
6. Що буде виведено за виконання фрагменту коду?

%code
print(5/5%5*5)
%code
\vspace{5mm}
%enditem
%beginitem
7. %goodpath:Scripts/2019/Mod2019_1/Lex/03-good.txt
Відмітити все, що є коректним оператором С++:
( 1 ) if a<b: a=0 else: b=0

( 2 ) x=3**1

( 3 ) if ('x'><'y'): a=6

( 4 ) --n
\vspace{5mm}
%enditem
%beginitem
8. Що буде виведено за виконання фрагменту коду?

%code
print('R', end=' ')
d=[10,11,12,13,14]
del d[-6:-6]
print(d)
%code

\vspace{5mm}
%enditem
%beginitem
9. 
Вкажіть значення та тип виразу:
not5.0>=3.0 and notTrue<=5 and 3!=7
\vspace{5mm}
%enditem
%beginitem
10. %goodpath:Scripts/2018/Mod2018_1/01-good.txt
Відмітити все, що є однією правильною лексемою:
( 1 ) 0,5

( 2 ) LIFO

( 3 ) x+23

( 4 ) x23
\vspace{5mm}
%enditem










%end
\newpage

\newpage

%begin
{\large\textbf{Контрольна робота}}

\textbf{Варіант \No\,9}\hfill \textit{ПІБ:}\,\underline{\hspace{60mm}}
\vspace{4mm}

%beginitem
1. Що буде виведено за виконання фрагменту коду?
%code
if (3 >= 3)
    cout << "x";
else
    cout << "b";
%code

\vspace{5mm}
%enditem
%beginitem
2. Що буде виведено за виконання фрагменту коду?
code
_d = 0

class D:
    _d = 1
    
    def __init__(self):
        self._d = 2
        _d = 3
        D._d = 4

obj = D()
try:
    print(_d, end=' ')
    print(obj._d, end=' ')
    print(D._d, end=' ')
    print(obj._d, end=' ')
    print(obj._D__d, end=' ')
except:
    print('ERR')
code
\vspace{5mm}
%enditem
%beginitem
3. %goodpath:Scripts/2019/Mod2019_1/Lex/03-good.txt
Відмітити все, що є коректним оператором С++:
( 1 ) if a<b: a=0 else: b=0

( 2 ) x=3**1

( 3 ) if ('x'><'y'): a=6

( 4 ) --n
\vspace{5mm}
%enditem
%beginitem
4. %goodpath:Scripts/2018/Mod2018_1/03-good.txt
Відмітити все, що є коректним оператором С++:
( 1 ) cout<<(3=!5);

( 2 ) bool f(int a);

( 3 ) cin<<x;

( 4 ) j=i;
\vspace{5mm}
%enditem
%beginitem
5. Що буде виведено за виконання фрагменту коду?

%code
print('R', end=' ')
d=[10,11,12,13,14]
del d[-6:-6]
print(d)
%code

\vspace{5mm}
%enditem
%beginitem
6. Що буде виведено за виконання фрагменту коду?

%code
def h(a,b):
    c=69
    if b>=-3:
        return 3
    if a<=-4:
         c=1
    else: 
        return 5
    return c

print(h(9,9))
%code
\vspace{5mm}
%enditem
%beginitem
7. Що буде виведено за виконання фрагменту коду?

%code
n=3
while n<7:
    n+=1
print(n)
%code

\vspace{5mm}
%enditem
%beginitem
8. %goodpath:Scripts/2019/Mod2019_1/Lex/01-good.txt
Відмітити все, що є однією правильною лексемою:
( 1 ) (1)

( 2 ) 0.2E3

( 3 ) "abc'

( 4 ) x23
\vspace{5mm}
%enditem
%beginitem
9. Що буде виведено за виконання фрагменту коду?

%code
print('R', end=' ')
n=7
while n>3:
    n+=-3
print(n)
%code

\vspace{5mm}
%enditem
%beginitem
10. Що буде виведено за виконання фрагменту коду?

%code
print('R', end=' ')
n=3
while n<7:
    n+=1
print(n)
%code

\vspace{5mm}
%enditem










%end
\newpage

\newpage

%begin
{\large\textbf{Контрольна робота}}

\textbf{Варіант \No\,10}\hfill \textit{ПІБ:}\,\underline{\hspace{60mm}}
\vspace{4mm}

%beginitem
1. Що буде виведено за виконання фрагменту коду?
d=[10,11,12,13,14]; del d[-6:-6]; print(d)\vspace{5mm}
%enditem
%beginitem
2. 

\begin{center}
	\adjustimage{max size={0.8\linewidth}{0.8\paperheight}}
	{../BinTree/trees_exp/59.png}
\end{center}
Указати послідовність міток вершин дерева, зображеного на рис., що утвориться в результаті обходу дерева в прямому порядку. Мітки записувати через пробіл.\par Обхід дерева починається з кореня (на рис. це верхня вершина).
%answer:59 прямому
\vspace{5mm}
%enditem
%beginitem
3. %goodpath:Scripts/2018/Mod2018_1/01-good.txt
Відмітити все, що є однією правильною лексемою:
( 1 ) 0,5

( 2 ) LIFO

( 3 ) x+23

( 4 ) x23
\vspace{5mm}
%enditem
%beginitem
4. Що буде виведено за виконання фрагменту коду?
try:
    print(3, end=';')
    print(1j>=3j, end=';')
    print(5, end=';')
except TypeError:
    print(0, end=';')
except ValueError:
    print(9, end=';')
except:
    print('Z', end=';')
else:
    print(6, end=';')
finally:
    print(4)
\vspace{5mm}
%enditem
%beginitem
5. %goodpath:Scripts/2018/Mod2018_1/03-good.txt
Відмітити все, що є коректним оператором С++:
( 1 ) cout<<(3=!5);

( 2 ) bool f(int a);

( 3 ) cin<<x;

( 4 ) j=i;
\vspace{5mm}
%enditem
%beginitem
6. Що буде виведено за виконання фрагменту коду?

%code
print(5/5%5*5)
%code
\vspace{5mm}
%enditem
%beginitem
7. %goodpath:Scripts/2019/Mod2019_1/Lex/03-good.txt
Відмітити все, що є коректним оператором С++:
( 1 ) if a<b: a=0 else: b=0

( 2 ) x=3**1

( 3 ) if ('x'><'y'): a=6

( 4 ) --n
\vspace{5mm}
%enditem
%beginitem
8. Що буде виведено за виконання фрагменту коду?

%code
#include <iostream>

int h(int a, int b){
    int c = 69;
    if (b >= -3)
        return 3;
    if (a <= -4)
         c = 1;
    else 
        return 5;
    return c;
}

int main(){
    std::cout << h(9, 9);
    return 0;
}
%code
\vspace{5mm}
%enditem
%beginitem
9. Що буде виведено за виконання фрагменту коду?
try:
    res = "3">=3.0
    if isinstance(res, bool):
        print(f'bool;{res}')
    elif isinstance(res, int):
        print(f'int;{res}')
    elif isinstance(res, float):
        print(f'float;{res:.2f}')
    else:
        print('???')
except:
    print('Z')
\vspace{5mm}
%enditem
%beginitem
10. Що буде виведено за виконання фрагменту коду?
def f(n):
    print('f', end='')
    return n

try:
    print(f(9) and f(-8))
except:
    print('ERR')
\vspace{5mm}
%enditem










%end
\newpage
\end{document}


